\chapter{Literature Review / Related Work}

\section{Definitions, Acronyms and Abbreviations}

\subsection{Definitions}
\begin{description}
    \item[Evapotranspiration] The combined process of water loss from the earth's surface into the atmosphere through evaporation from soil and transpiration from plants.
    \item[Reference Evapotranspiration ($ET_0$)] The evapotranspiration rate from a well-watered reference surface, typically an extensive surface of green grass of uniform height.
    \item[Crop Coefficient ($K_c$)] A dimensionless ratio that relates the reference evapotranspiration to the actual crop evapotranspiration under standard conditions.
    \item[Mandi] A local or regional wholesale agricultural market where farmers sell their produce to commission agents or buyers.
    \item[Time Series Forecasting] The use of mathematical and statistical models to predict future values based on previously observed, chronologically ordered data.
    \item[Multi-Criteria Decision Making (MCDM)] A sub-discipline of operations research that explicitly evaluates multiple, often conflicting, criteria in decision making.
    \item[Soil Moisture] The quantity of water contained in the soil matrix, which is a critical variable for determining irrigation triggers.
    \item[Precision Agriculture] A farming management concept based on observing, measuring, and responding to inter- and intra-field variability in crops.
    \item[Vernacular Interface] A user interface localized into a native or indigenous language, tailored directly to users who do not speak English.
\end{description}

\subsection{Acronyms and Abbreviations}
\begin{description}
    \item[ANN] Artificial Neural Network
    \item[ARIMA] AutoRegressive Integrated Moving Average
    \item[CNN] Convolutional Neural Network
    \item[DSS] Decision Support System
    \item[ET\textsubscript{0}] Reference Evapotranspiration
    \item[FAO] Food and Agriculture Organization
    \item[GIS] Geographic Information System
    \item[LSTM] Long Short-Term Memory
    \item[MAE] Mean Absolute Error
    \item[MAPE] Mean Absolute Percentage Error
    \item[MCDM] Multi-Criteria Decision Making
    \item[ML] Machine Learning
    \item[NASA POWER] NASA Prediction of Worldwide Energy Resources
    \item[PAMRA] Punjab Agriculture Marketing Regulatory Authority
    \item[RMSE] Root Mean Square Error
    \item[UI] User Interface
    \item[UX] User Experience
\end{description}

\section{Detailed Literature Review}

\subsection{Agricultural Price Forecasting}

\subsubsection{Summary of the research items}
Menculini et al. \cite{menculini2021comparing} propose comparing deep learning methods such as Prophet alongside traditional ARIMA architectures for forecasting wholesale food prices. Using a large dataset of regional wholesale market prices, the study demonstrates that deep learning significantly outperforms classical statistical techniques, achieving lower error rates across different food commodity models.

In \cite{murugesan2022forecasting}, Murugesan and Mishra evaluate various deep learning architectures, including basic LSTM, bi-LSTM, stacked LSTM, and CNN-LSTM configurations, for predicting agricultural commodity prices. By extracting temporal sequences and spatial patterns, this study demonstrates that CNN-LSTM hybrids offer a highly accurate forecasting mechanism, outperforming basic recurrent neural networks in complex market scenarios.

The study by Joshi and Patel \cite{joshi2022cnn} utilized a CNN-Bidirectional LSTM approach specifically applied to price forecasting of volatile agricultural commodities in Gujarat. The bidirectional model effectively learned complex non-linear trading patterns, confirming the suitability of deep sequential models to capture both forward and backward temporal dependencies in local agriculture markets.

Finally, Habib et al. \cite{habib2021price} analyse the impact of price volatility on agricultural commodities versus general food items with a specific focus on Pakistan. Utilizing localized economic datasets alongside historical prices, the authors demonstrate how macroeconomic instability and international commodity trends compound short-term price fluctuations. Their baseline findings underscore the necessity for robust, localized predictive tools to mitigate these shocks.

\subsubsection{Critical analysis of the research items}
The comparative study by Menculini et al. \cite{menculini2021comparing} successfully establishes the supremacy of deep learning over ARIMA; however, its reliance on stable European datasets limits its direct generalisability to highly erratic markets like Pakistan's mandi system. The deep learning models evaluated by Murugesan and Mishra \cite{murugesan2022forecasting} are notably powerful, yet complex architectures such as CNN-LSTM run the risk of overfitting when applied to data-scarce environments. The bidirectional approach from Joshi and Patel \cite{joshi2022cnn} handles long-term dependencies gracefully but lacks robust handling of exogenous variables, such as sudden policy shifts or extreme weather shocks, which regularly distort agricultural pricing. Meanwhile, Habib et al. \cite{habib2021price} correctly identify structural volatility issues in Pakistan's agriculture sector, but the work remains heavily diagnostic rather than predictive, offering no implementable machine learning framework for regional stakeholders.

\subsubsection{Relationship to the proposed research work}
The CNN-LSTM architecture evaluated positively by Murugesan and Mishra \cite{murugesan2022forecasting} translates to the baseline design for KisaanMadad's price forecasting module, effectively balancing spatial feature extraction with time-series memory similar to the approach by Joshi and Patel \cite{joshi2022cnn}. To address the localized volatility highlighted by Habib et al. \cite{habib2021price}, the module will rely on PAMRA and AMIS daily mandi price datasets spanning 2018--2023, rather than abstract global commodity curves. This bridges the gap between state-of-the-art deep learning findings and the resource-constrained needs of smallholder farmers in Pakistan's localized supply chains.

\subsubsection{Literature review summary table}
\begin{table}[htbp]
  \centering
  \caption{Literature review summary table for Agricultural Price Forecasting}
  \label{tab:price_forecasting}
  \begin{tabular}{p{0.15\linewidth} p{0.25\linewidth} p{0.25\linewidth} p{0.25\linewidth}}
    \toprule
    \textbf{Authors} & \textbf{Method} & \textbf{Results} & \textbf{Limitations} \\
    \midrule
    Murugesan \& Mishra \cite{murugesan2022forecasting} & CNN-LSTM and LSTM variants & Superiority of CNN-LSTM architectures & High computational cost, data heavy \\
    Joshi \& Patel \cite{joshi2022cnn} & CNN-Bidirectional LSTM & Accurate regional forecasting in Gujarat & Ignores exogenous market shocks \\
    Menculini et al. \cite{menculini2021comparing} & Prophet, Deep, ARIMA & Prophet and DL outperformed ARIMA & Tested on stable European datasets \\
    Habib et al. \cite{habib2021price} & Economic Analysis & Identified severe structural price volatility & Diagnostic only, no predictive application \\
    \bottomrule
  \end{tabular}
\end{table}

\subsection{Crop Recommendation Systems}

\subsubsection{Summary of the research items}
Raghavendra et al. \cite{raghavendra2024systematic} present an extensive systematic literature review of machine learning algorithms utilized strictly for crop recommendation and yield estimation. Evaluating a consolidated body of work, the study demonstrates that ensemble methods and deep learning models consistently outperform traditional ML approaches when trained on diverse soil and atmospheric inputs.

In \cite{priyanka2023iot}, Priyanka et al. present an IoT-based crop recommendation framework using a machine learning backend for smart agriculture. Using sensory data measuring temperature and moisture, this study maps real-time environmental data to crop viability, demonstrating early success in automating recommendations dynamically.

Kancharagunta et al. \cite{kancharagunta2024crop} survey crop suitability systems by combining diverse IoT inputs with machine learning models. The survey highlights robust methodologies that integrate continuous environmental parameters into algorithmic decisions, successfully aiding both yield prediction and agricultural scheduling.

Finally, Apat et al. \cite{apat2023artificial} focuses specifically on artificial intelligence-based frameworks for crop guidance, proposing a lightweight intelligent recommendation system. By classifying secondary datasets encompassing soil and water availability features, the proposed system proves highly capable for farmers lacking expensive localized hardware.

\subsubsection{Critical analysis of the research items}
While Raghavendra et al. \cite{raghavendra2024systematic} provides an excellent foundational synthesis of machine learning in agriculture, the aggregated literature flattens regional heterogeneities, making perfect accuracy difficult to replicate in real-world scenarios. The IoT-driven architecture by Priyanka et al. \cite{priyanka2023iot} succeeds functionally but remains entirely cost-prohibitive for the mass majority of farmers owning under 12.5 acres. The survey approach by Kancharagunta et al. \cite{kancharagunta2024crop} confirms these hardware-adoption barriers. Conversely, the methodology by Apat et al. \cite{apat2023artificial} is a contextually robust software-only approach, yet its reliance on secondary static traces typically limits its predictive power against volatile precipitation cycles and climate changes.

\subsubsection{Relationship to the proposed research work}
KisaanMadad's crop recommendation system builds upon the software-centric philosophy presented by Apat et al. \cite{apat2023artificial} while aiming for the robust accuracy highlighted by Raghavendra et al. \cite{raghavendra2024systematic}. To overcome the lack of IoT hardware noted as a vulnerability by Priyanka et al. \cite{priyanka2023iot} and Kancharagunta et al. \cite{kancharagunta2024crop}, KisaanMadad will integrate dynamic, real-time meteorological inputs via the NASA POWER API alongside FAO crop calendars. By weighing continuous remote variables against static local criteria, KisaanMadad proposes an accurate recommendation system directly tailored to small-scale growers in Punjab.

\subsubsection{Literature review summary table}
\begin{table}[htbp]
  \centering
  \caption{Literature review summary table for Crop Recommendation Systems}
  \label{tab:crop_recommendation}
  \begin{tabular}{p{0.15\linewidth} p{0.25\linewidth} p{0.25\linewidth} p{0.25\linewidth}}
    \toprule
    \textbf{Authors} & \textbf{Method} & \textbf{Results} & \textbf{Limitations} \\
    \midrule
    Priyanka et al. \cite{priyanka2023iot} & IoT-based ML classification & High dynamic accuracy & Hardware-dependent, cost-prohibitive for smallholders \\
    Apat et al. \cite{apat2023artificial} & AI modeling on secondary data & Highly viable crop mapping & Relies heavily on static or prior weather traces \\
    Kancharagunta et al. \cite{kancharagunta2024crop} & Survey of ML/IoT systems & Synthesis of precision agritech adoption & Implementation complexity highlighted in survey \\
    Raghavendra et al. \cite{raghavendra2024systematic} & Systematic literature review & Identified optimal ensemble models & Generalised findings ignore regional variances \\
    \bottomrule
  \end{tabular}
\end{table}

\subsection{FAO-56 Irrigation Scheduling and ET\textsubscript{0} Estimation}

\subsubsection{Summary of the research items}
The foundational work by \cite{allen1998crop} establishes the global standard for computing crop water requirements via the FAO-56 Penman-Monteith method. This exhaustive procedural framework uses defined crop coefficients ($K_c$) and standard meteorological variables to produce the definitive reference evapotranspiration ($ET_0$) metric used globally.

In \cite{ferreira2020smartphone}, the authors propose a smartphone application specifically architected for weather-based irrigation scheduling through Artificial Neural Networks. By processing atmospheric variables against known evapotranspiration baselines in a localized dataset, the app efficiently minimized irrigation errors, demonstrating an inherently low RMSE of 0.82 mm/day compared to manual timing approximations.

A study by Baber and Ullah \cite{baber2024short} evaluates the temperature-based Hargreaves-Samani equation against the standard FAO-56 Penman-Monteith equation for $ET_0$ estimation strictly within Pakistan. Using regional observatory data, this study demonstrates that while the PM method holds theoretical supremacy, the calibrated Hargreaves-Samani model is exceptionally suitable and accurate for data-scarce Pakistani regions.

Lastly, Jaafar et al. \cite{jaafar2022agsat} proposes AgSAT, a smart irrigation application reliant on cloud-fetched NASA POWER data combined with satellite imagery. By directly calculating daily crop ET and water requirements via remote methodologies, the authors demonstrate robust irrigation scheduling capabilities without physical ground equipment.

\subsubsection{Critical analysis of the research items}
The canonical methodology provided by \cite{allen1998crop} remains technically unmatched, yet its rigid demand for numerous meteorological variables makes it exceptionally difficult to directly deploy in sensor-less regions. The smartphone interface of \cite{ferreira2020smartphone} democratizes irrigation insights but requires continuous real-time connectivity to process neural network computations, structurally isolating off-grid users. The regional evaluation by Baber and Ullah \cite{baber2024short} highlights a crucial compromise: the Hargreaves equation is simpler, but without proper calibration coefficients, farmers risk wasteful over-irrigation. The utilization of NASA POWER datasets by mobile applications like the one developed by Jaafar et al. \cite{jaafar2022agsat} presents an optimal software pathway, though resolution constraints inevitably restrict effectiveness on extreme micro-topographies.

\subsubsection{Relationship to the proposed research work}
Following the data sourcing successes of Harris \cite{harris2025energy} and Jaafar et al. \cite{jaafar2022agsat}, KisaanMadad bypasses physical hardware constraints by querying the NASA POWER API for real-time solar radiation and temperature. This remote data populates the core FAO-56 Penman-Monteith engine \cite{allen1998crop} compiled within the application. When wind speed or relative humidity inputs fail, KisaanMadad falls back to the temperature-driven Hargreaves methodology verified for use in local contexts by Baber and Ullah \cite{baber2024short}, guaranteeing uninterrupted irrigation scheduling.

\subsubsection{Literature review summary table}
\begin{table}[htbp]
  \centering
  \caption{Literature review summary table for FAO-56 Irrigation Scheduling}
  \label{tab:irrigation}
  \begin{tabular}{p{0.15\linewidth} p{0.25\linewidth} p{0.25\linewidth} p{0.25\linewidth}}
    \toprule
    \textbf{Authors} & \textbf{Method} & \textbf{Results} & \textbf{Limitations} \\
    \midrule
    Harris \cite{harris2025energy} & Edge AI and NASA POWER & Minimized energy load and water use & Requires continuous edge/cloud connectivity \\
    Baber \& Ullah \cite{baber2024short} & Penman-Monteith vs Hargreaves & Verified Hargreaves-Samani suitability & Calibration issues without proper coefficient setup \\
    Jaafar et al. \cite{jaafar2022agsat} & NASA POWER \& Satellite apps & Strong field-scale ET metrics & Satellite resolutions may overlook micro-climates \\
    Allen et al. \cite{allen1998crop} & FAO-56 foundational standard & Global ET0 benchmark established & Requires extensive meteorological input variables \\
    \bottomrule
  \end{tabular}
\end{table}

\subsection{Agricultural Mobile Applications in Developing Countries}

\subsubsection{Summary of the research items}
Kamal and Bablu \cite{kamal2023mobile} propose a systematic review regarding the widespread adoption of mobile applications by smallholder farmers operating in developing nations. Their analysis highlights how applications empowering farmers can accelerate agricultural development, provided digital payment systems and information delivery align with ground truths.

In \cite{amina2025factors}, Amina et al. specifically evaluate the factors influencing farmers' satisfaction and adoption of mobile phone applications for agricultural purposes within Faisalabad, Pakistan. The study found that usability and reliability directly determine satisfaction, pushing young farmers to optimize operations through reliable, well-designed mobile interfaces.

The investigation by Shahzad et al. \cite{shahzad2024mobile} examines digital agriculture adoption across Punjab, Pakistan, assessing the role of perceived awareness on farmer decisions. Utilizing a large random sample from diverse agroecological zones, the study proves that localized, culturally coherent outreach dramatically enhances platform adoption and overall digital proficiency.

Finally, examining existing Pakistani platforms, the \cite{pasha2024agritech} report reviews the state of fragmented agricultural tech solutions. Summarizing applications capable of either market updates or weather warnings, the report demonstrates an absolute failure among developers to integrate diverse utilities into unified systems, forcing farmers to operate multiple disparate applications for daily problem-solving.

\subsubsection{Critical analysis of the research items}
The review by Kamal and Bablu \cite{kamal2023mobile} comprehensively diagnoses the positive macro-impacts but leaves technical implementation details, like offline caching or UI constraints, unaddressed. The regional usability analysis by Amina et al. \cite{amina2025factors} correctly identifies young farmers as the initial adoption catalyst in Pakistan, but overlooks the usability constraints faced by aging or low-literacy demographics. The Punjab-centric study by Shahzad et al. \cite{shahzad2024mobile} provides localized adoption data entirely relevant to KisaanMadad, yet highlights awareness, rather than purely technical features, as the core issue. Furthermore, the Pasha \cite{pasha2024agritech} sector report critically highlights app fragmentation, emphasizing that single-feature platforms fail to retain user interest without a comprehensive ecosystem backend.

\subsubsection{Relationship to the proposed research work}
Addressing the exact usability hurdles documented by Amina et al. \cite{amina2025factors} and the necessity of targeted local demographic awareness proven by Shahzad et al. \cite{shahzad2024mobile}, KisaanMadad is natively engineered with a vernacular Urdu User Interface utilizing accessible, heavily iconized navigational paradigms. Responding directly to the crippling fragmentation decried by the Pasha \cite{pasha2024agritech} report and following the empowerment guidelines set by Kamal and Bablu \cite{kamal2023mobile}, KisaanMadad consolidates price forecasting, crop recommendation, irrigation scheduling, and financial planning into one unified platform, representing a holistic, unprecedented counter to the limitations of current rural deployments.

\subsubsection{Literature review summary table}
\begin{table}[htbp]
  \centering
  \caption{Literature review summary table for Agricultural Mobile Applications}
  \label{tab:mobile_apps}
  \begin{tabular}{p{0.15\linewidth} p{0.25\linewidth} p{0.25\linewidth} p{0.25\linewidth}}
    \toprule
    \textbf{Authors} & \textbf{Method} & \textbf{Results} & \textbf{Limitations} \\
    \midrule
    Kamal \& Bablu \cite{kamal2023mobile} & Review of agricultural dev apps & Proven tech empowerment linkages & Diagnoses macro trends without UI technical guidance \\
    Amina et al. \cite{amina2025factors} & Faisalabad adoption survey & Identified reliability as key driver & Predominantly evaluates younger demographics \\
    Shahzad et al. \cite{shahzad2024mobile} & Punjab field adoption stats & Strong correlation with local awareness & Awareness-centric, lacks focus on app mechanics \\
    Pasha \cite{pasha2024agritech} & Industry sector reporting & Identified severe app fragmentation & Enterprise focus, no underlying predictive models \\
    \bottomrule
  \end{tabular}
\end{table}

\subsection{Related Applications}

\subsubsection{Summary of the research items}
The Kisan App \cite{kisanapp2024} is an established digital advisory platform targeting farmers primarily within Punjab, Pakistan. It focuses heavily on static farming methodologies, regional news alerts, and general weather reporting. It boasts a large user base due to widespread government and telco collaboration, effectively serving as an elementary portal for regional agricultural information.

Farmdar \cite{farmdar2024} is an established agritech business offering a suite of B2B products (such as CropScan and Yield Pro) rather than a single farmer-facing application. It employs high-resolution, multi-spectral remote sensing AI over satellite imagery to monitor crop health at scale. Deployed heavily across sugar mills, corporate farms, and lending institutions throughout Pakistan and globally, it specializes in actionable spatial detection regarding nutrient deficiencies, pest outbreaks, and yield predictions, driven directly by precision GIS analytics.

BaKhabar Kissan \cite{bakhabarkissan2024} is undoubtedly Pakistan's most prominent holistic agricultural advisory platform. Supported by a robust backend call-center infrastructure and an aggressive telecommunications partnership, the application delivers localized weather parameters, pest warnings, and up-to-date regional market commodity prices. It caters extensively to smallholder farmers seeking daily operational guidance.

\subsubsection{Critical analysis of the research items}
While the Kisan App \cite{kisanapp2024} is accessible, its feature set is profoundly rudimentary; it severely lacks any advanced machine learning integration for intelligent crop recommendations, and it does not possess price forecasting or financial calculators, rendering it purely informational rather than a true decision support system. Farmdar \cite{farmdar2024} leverages state-of-the-art spatial precision, yet its enterprise focus, B2B product structure, and heavy reliance on extensive land tracts organically alienate the 97\% of farmers confined to holdings under 12.5 acres. Conversely, BaKhabar Kissan \cite{bakhabarkissan2024} addresses smallholders effectively through price reporting, but critically acts only as a diagnostic repository; it fails entirely to forecast incoming price volatility, offers no sophisticated FAO-56 integrated localized irrigation scheduling, and completely omits modular financial planning.

\subsubsection{Relationship to the proposed research work}
The gap analysis of existing entities inherently cements KisaanMadad's unique technological positioning. While BaKhabar Kissan \cite{bakhabarkissan2024} provides basic prices, KisaanMadad introduces dynamic models predicting future prices \cite{murugesan2022forecasting} previously unavailable at zero cost. Unlike the enterprise-focused Farmdar products \cite{farmdar2024}, KisaanMadad explicitly targets smallholders by replacing cost-prohibitive remote B2B analytics with multi-criteria recommendations and FAO-56 \cite{allen1998crop} frameworks operating natively on mobile. Finally, surpassing the passive information structures of the Kisan App \cite{kisanapp2024}, KisaanMadad is the first architecture to fuse agronomic tools with foundational financial planning resources, closing the loop between ecological farm management and smallholder macroeconomic survival.

\subsubsection{Related applications summary table}
\begin{table}[htbp]
  \centering
  \caption{Related applications summary table}
  \label{tab:related_apps}
  \begin{tabular}{p{0.15\linewidth} p{0.30\linewidth} p{0.30\linewidth} p{0.15\linewidth}}
    \toprule
    \textbf{Application} & \textbf{Features} & \textbf{Relevance to KisaanMadad} & \textbf{Limitations} \\
    \midrule
    Kisan App \cite{kisanapp2024} & News, static advice, weather & General template for UI accessibility & No ML, no forecasting, no financial tools \\
    Farmdar \cite{farmdar2024} & Remote sensing B2B products & Demonstrates success of precision agritech & Expensive, enterprise-focused, not a direct consumer app \\
    BaKhabar Kissan \cite{bakhabarkissan2024} & Daily prices, pest warnings, call centre & Key competitor in price dissemination & Passive reporting only, no FAO-56 irrigation \\
    \bottomrule
  \end{tabular}
\end{table}
% === END CHAPTER 3 ===
